\chapter{Why Vote for School Board?}
\label{chap:turnout}
\chapterstyle{deposit}




\section{Introduction}
Since \citet{Downs57} political scientists have been asking the question: why vote? 
Voters in local school board elections may have been asking this question even 
earlier. Despite the fact that in general school boards are in charge of levying 
a substantial share of many citizens' tax bill through property taxes and 
collectively manage expenditures annually equivalent to the budget of the US 
Department of Defense, when school board elections are held off-cycle turnout 
and voter interest in them appears to be remarkably low \citep{RefWorks:333, RefWorks:321}. 
This despite the fact that the community stake in quality education is large, not 
just for parents of the students, but for the social, economic, and political 
impacts that high or low quality schools can have on a community. %citation needed

We know that voter turnout even in national on-cycle elections in the US remains 
lower than many other democracies \citep{BlaisTurnout2006}. We also have an extensive 
literature seeking to explain patterns in turnout between years and within years 
between social groups \citep{ParticAnnReview2004, econvoting2007, Galston2001, EconandOutcomes2000}.
The specific puzzle of voter turnout in school board elections has even spawned 
its own set of theories with empirical tests \citep{RefWorks:311,RefWorks:243,RefWorks:340}.

Despite this rich ground, much work remains to be done to adequately test the 
mechanisms at play in predicting turnout in off-cycle, local, non-partisan elections 
like school board elections. Furthermore, theories of participation in on-cycle and 
higher profile elections benefit from examining the limits of their generalizability and 
exploring the mechanisms that contribute to this. In this paper I seek to explore 
voter turnout in school board elections through the lenses of the major theories of 
voter turnout. I begin by reviewing these theories. Then, I move into a discussion 
of school boards and the contribution a study of school boards can bring to 
the understanding of these theories. Then I discuss the unique panel data set of 
school board results that I have collected and the methods with which I will test 
these competing theories on the data. Finally, I conclude with an assessment of 
the performance of each of the theories and the foundations for future work. 

%% Significance of school boards
%% - as a share of elections they are dominant
%% - they have tremendous consequences to parents and tax payers 
%% - they are very different structurally, more variation, richer ground for theory


\section{Why Vote?}

Setting aside important though largely settled issues of access to the polls, 
enfranchisement laws, administrative and cognitive barriers to casting a ballot, and 
freedom from intimidation -- all thankfully of minimal concern in American school 
board elections -- there remains a wealth of intricate and interrelated factors 
which shape a voter's decision to cast a ballot.%CITATION needed here
I broadly group this literature 
into three major strands and include the theories that have sprung up uniquely 
around school boards as a fourth category. My aim is not to fully summarize this 
rich theoretical ground, but to reframe these debates in light of the particular 
challenge of school board elections. As prologue, I succinctly summarize the 
unique case against voting in school board elections as follows:

\begin{itemize}
  \item I don't know when the election is, or I forgot. (salience)
	\item I don't know who is running. (attention)
	\item I don't have a choice - no one is challenging the incumbent. (opportunity)
	\item I don't care which candidate wins. (lack of partisan cues)
	\item I'm satisfied with my schools. (lack of impetus)
	\item I don't care how schools perform or what the property tax is (salience)
	\item My vote won't matter
\end{itemize}

These concerns of our average potential voter are important to keep in mind as 
the major strands of voter participation theory are sketched out in the next 
sections. Other researchers studying municipal or provincial elections have 
found similar concerns with turnout \citep{Hajnal01052003, PPIC2002}. 


The turnout literature asks two fundamental questions -- why does 
anyone vote in the first place, and why do some people vote sometimes, and others 
do not? This chapter will attempt to answer these questions for off-cycle 
school board elections in Wisconsin. 

\begin{itemize}
	\item Rational Choice
	\item Group Mobilization
	\item Information
  \item Habits and Socialization
\end{itemize}



\subsection{Rational Choice}

\citet{Downs57}'s elegant statement of the problem of voting as a question of 
why anyone would vote when their vote has such a small probability of influencing 
the outcome presented political scientists with a challenge of understanding why 
people do vote. The now famous formal statement of the problem as given by Downs:

\[ R = PB - C > 0 \]

R is the net expected utility of voting, P is the probability of influencing the 
election, B is the difference in the expected utility of the policies of the 
two candidates, and C refers to the cost of voting. \citet{VoteCalculus68}'s 
refined the equation further by adding the D term:

\[ R = PB - C + D > 0 \]

which was intended to capture the benefit of self-expression associated with voting -- 
either self-expression of candidate preference, or self-expression of the utility of 
voting. However, the D-term merely shifted the debate - the probability of benefit 
from an individual vote still remained zero, so voters would only vote when $D > C$ - 
that is, when the benefit of self-expression was greater than the cost of voting 
\citep{nonvoting74, feddersenreview}.\footnote{\citet{Dowding2005} places an interesting 
wrinkle on this by questioning why voters care about being the decisive vote when 
politicians so often seek large victories to ensure a safe seat or to earn a mandate. 
This seems less germane to the present case - school board members and the community 
are unlikely to interpret a substantial victory as a mandate.} 
While the D term did not solve the puzzle posed by \citet{Downs57}, it 
did touch off a decades long research project to uncover why the D term varies 
among voters, and the functional form of this benefit of self-expression.\citep[see:][]{KanYang2001, 
Dowding2005,feddersenreview,Schuessler1} 

%%%%%%%%%%% CITES %%%%%%%%%%%%%%%%%

While I will not attempt to summarize this theoretically rich debate here 
\citep[see instead:][]{feddersenreview,geysreview}, I do want to highlight that most empirical 
tests of these voter turnout theories have focused on large-scale electorates such 
as state or national elections. There exists a class of voter turnout theories that 
predict very different behaviors for differently sized electorates. School board 
elections are typically comprised of small electorates (STATISTIC NEEDED) where 
P is much higher than other elections.\footnote{P can be calculated following the 
method of (CITATION) for school board races.} It remains to be seen if B is indeed 
higher in such cases, as we shall see the policy preferences of candidates for 
board are not always easily discovered by voters. Yet, game theoretic and social 
mobilization theories of turnout expect that in smaller electorates the collective 
behavior of the electorate or social group could spur turnout to be higher. 

% Game theoretic

In the game theoretic case the argument is that if everyone pursues the optimal 
strategy of not voting, then the value of PB increase. This sets of a strategic 
game among voters where the value of P is determined by the strategies selected by 
all other voters in the electorate \citep{Palfrey83, Palfrey85,Ledyard84}. In order for positive turnout to emerge from 
such a game, voters are expected to have certainty about the costs of voting and 
the preferences of other voters -- an unrealistic assumption in large electorates that 
hurts the case for such an approach \citep{Mueller3, aldrich1993}. This approach 
is often called the pivotal-voter model, and it has been explicitly tested in 
small-scale elections and found that while it predicts overall turnout sufficiently, 
it does not give an accurate reflection of the margins of victory \citep{Coate2008582}. 
Other scholars find size to be an important factor in local elections, but do not 
tie this to a particular theory of voter turnout \citep{Frandsen2002}.\footnote{The size 
question is a curious one, because at its extreme it would predict that lower-level 
elections should thus exhibit higher turnout than national elections - a phenomenon 
described by \citet{Horiuchi2001} as the ``turnout-twist.''}

%%%%%%%%%%% CITES %%%%%%%%%%%%%%%%%

\subsection{Group Mobilization}

Social networks and group organization might also help explain voter turnout. \citet[p.23]{geysreview}
states the key of the group-based models:

\begin{quote}
First, groups are likely to have larger benefits than individuals from political participation. 
The reason is that politicians may provide groups with extra benefits – in terms of 
policies that come closer to the group’s optimum – to win the support of the group \citep{Lapp99}. 
Second, as the political influence of a social group can be assumed to be proportional to its size
\citep{SchrWinden1991}, the group as a whole is more likely to have a 
non-negligible impact on the election outcome. 
\end{quote}

The group approach appears to more accurately reflect the modern political environment 
and campaign strategy of mobilization and targeted messaging to voters \citep{Lapp99}. 
The theory is more nuanced because it allows for the case where the optimal strategy 
for a group is to abstain from voting \citep{geysreview, Fowler2005}. 
Crucially, such group-based mobilization relies on the group being able to enforce 
the voting norm within the group. This enforcement has been shown to depend on 
factors that seem especially reflected within small-scale local elections. 
\citet[p.25]{GrossmanHelpman} identify three such factors: the 
frequency of interactions within the group, the risk of social isolation resulting 
from deviating from the group's behavior, and the ability of group members to 
monitor one another. These factors are essential to preventing free riding \citep[see][]{Olson65}.  
\citet{SchrWinden1991} identifies opinion leaders within these groups as a key 
transmitter of social pressure. These leaders have incentives to increase turnout such 
as increasing their influence with elected officianls, and voters wanting to build 
credibility with these leaders find that turning out to vote is important. In the case 
of school board elections such local effects might be amplified in communities with 
strong local organizations where monitoring costs are low. \footnote{\citet{Morton91} finds that there is an equilibrium 
with positive turnout if at least one leader has a strict preference for a single candidate.}

Notably missing from this summary is the influence that candidates and campaigns 
have on voter turnout. While there is some evidence that education messages used 
by candidates can be a mobilizing force for some voters and of the role campaign 
expenditures in local elections can play in mobilizing voters \citep{SidesKarch2008, 
Holbrook15072013}, 
the impact of candidate campaigning and voter mobilization in school board 
elections is likely marginal. In the most recent survey of school board members, over 
two-thirds of board members report their last election as being ``easy'' or ``somewhat 
easy'', and in small districts this percentage is even higher \citep{RefWorks:322}.
\footnote{What's more, it is hard to know how much mobilization or advertising 
candidates can do when nearly 75\% of survey respondents report spending less than 
\$1,000 on their last campaign, and in small districts over 95\% of candidates 
reported spending less than \$1,000 on their campaign.\citep{RefWorks:322}}

\subsection{The Role of Information}
An important addition to this theoretical base came by adding the influence of 
limited voter information to voter turnout \citep{Matsusaka}. Voters' uncertainty 
around the B term, the policy difference between candidates, makes voting a riskier 
decision even for voters with a high predisposition for voting. The better informed 
a voter is, the more certainty they have that their vote will increase the likelihood 
of their policy preferences being realized. Partisan cues play an important role 
in this framework by reducing the cost of acquiring certainty about policy preferences and
at times drawing clear distinctions between them \citep{aldrich1993}.\footnote{
\citet{Larcinese} adapted \citet{Matsusaka} by showing that non-partisan citizens are most 
likely to acquire information, and thereby increase their likelihood to vote. WHY?!?!} 

An important feature of this school of thought is that much of the empirical testing 
of these theories has been done in situations where voters have a high chance of being 
incidentally, that is without incurring costs, exposed to information about the election 
because of the high profile nature of these races \citep{aldrich1993}. The informational 
cues provided by partisan signals in such races even further reduces this. School board
races, for the most part, feature neither of these cost reducing features. It seems unlikely 
that these higher costs are accompanied by higher benefits in such elections. Though 
voters are assured a higher likelihood of casting the deciding vote, they are much less certain 
that the election of candidate A over candidate B will make much of a difference in 
their quality of life in the immediate or medium term. 

% School boards may prove to be the electoral unit for which the Downsian hypothesis 
% is the most strained. In many state school boards are sub-county administrative 
% entities without partisan office holders and who often elect multiple members in 
% one election cycle. Furthermore, the low level of information about school board 
% candidates, their positions on issues, their likely vote share in an election, and 
% the voters' own feelings on issues facing the school board make this environment 
% dramatically different than the presidential election. What's more, the likelihood 
% that the vote of any single voter matters in a school district is orders of 
% magnitude larger than for state wide or nation wide office. 

%%%%%%%%%%% CITES %%%%%%%%%%%%%%%%%

\subsection{Voter Socialization}

Other scholars have found that the D term can be endogenous for individuals from 
election to election \citep{Plutzer2002}. Experimental 
evidence has found that ``Voting and abstention, in other words, are habit forming'' 
\citep[p.540]{Gerber2003} This is related to the group mobilization arguments 
above because social groups may well be the mechanism by which the habit is 
formed \citep{Kanazawa2000}. Particularly in groups that interact frequently or 
in tight knit communities social rewards may be found for voting and sanctions 
may be made for those who abstain \citep{Fowler2005}. Then, citizens who are 
rewarded for their vote (via the election of their preferred candidate) or punished for their abstention
(through the election of a less-preferred candidate) acquire an increased prefer-
ence for voting. In contrast, if their voting is punished or their abstention is
rewarded, they lose some of their preference for voting. This has been considered 
in other cases by political scientists as well \citep{ParticAnnReview2004}.\footnote{
Other scholars have found that not just the depth and frequency of interactions with 
voters matters, but also the nature of those interactions, the ideological 
composition of those groups, and socio-economic status
\citep{Mutz2002, classpart04, classbias92, sesbias92}.}


In school boards there is evidence that such socialization has occurred in the case of 
conservative Christians \citep{RefWorks:240}, unionized teachers \citep[chapter 5]{RefWorks:249, RefWorks:207}, 
and suggestions about the role institutions may play in deepening or weakening 
such socialization among ethnic minorities \citep{RefWorks:195}. Evidence of socialization 
is occurring might be found by observing ``stickiness'' in voter turnout 
in school board elections year to year, independent of other factors. \citet{RefWorks:248} 
called this the social conditioning model of school board turnout and found it held 
explanatory power in California school board election voter registration rates, 
but not rates of voter turnout. If members of the community vote at a higher 
average level after a high turnout election, this may provide evidence that the 
influx of turnout in the prior election has influenced the current election. 

\subsection{Dissatisfaction Theory}

In Chapter \ref{chap:theoryreview} we introduced dissatisfaction theory as the 
dominant theoretical paradigm for understanding school board elections.
This theory traces its roots back to \citet{RefWorks:348}'s concept of critical 
elections and \citet{RefWorks:319}'s conception of pluralism. It describes an electoral system with relative stability and little 
involuntary incumbent turnover punctuated by periods of citizen 
dissatisfaction, contentious elections, and incumbent defeats \citep{RefWorks:297,RefWorks:271}. 
Figure \ref{dissatdiag} displays the sequence of events in the model. Missing from 
empirical tests of dissatisfaction theory to date are indicators of voter participation. 
Implied throughout the dissatisfaction theory literature is that incumbent defeat 
is accompanied by an upwelling in citizen dissatisfaction and presumably voter 
participation in the election. Incumbents are defeated because new voters enter 
the game. Yet, until now, voter turnout in school board elections has not been 
used to test dissatisfaction theory explicitly, though it has received limited 
attention in understanding other phenomena of local education politics \citep{RefWorks:248, 
RefWorks:249, RefWorks:207, RefWorks:218}. \footnote{Taxpayers also have another 
option for expressing their dissatisfaction, which is to move jurisdictions \citep{MinkoffLyons}. 
In Wisconsin, yet another option of expressing dissatisfaction with the school board 
is available to parents - open-enrolling out of the school district. Both of these 
forms of dissatisfaction are decidedly higher cost than voting in a school board 
election, though their benefit is more immediate and obvious.}


Low-salience and low-participation elections thus can be democratic, if there is 
evidence that they are punctuated by periods of electoral activity when needed. 



\section{School Boards and Voter Turnout}

The current project adds importantly to the puzzle of voter turnout in two important ways. 
First, it provides an important update to the prior literature on school board 
election participation by explicitly measuring voter turnout in school board 
elections and investigating several hypothesized drivers of voter participation 
in such elections \citep{RefWorks:311,RefWorks:243,RefWorks:340}. Second, this study 
explores the linkage between state level politics and local electoral participation. 
Thus, the study represents an attempt to estimate the causal effect of a state 
level policy shock on voter turnout in local elections by leveraging the unique 
conditions in Wisconsin from 2010-2012. The present study presents an important 
oppotunity to test theories of participation and cross-level linkages in a federal 
system that previously were not available to political scientists. 

\subsection{Voter Turnout in School Board Elections}

While some rational choice models have been adapted to the school board contest, 
there exist no large scale studies of voter turnout in school board elections 
that have empirically tested these theories 
\citep{RefWorks:243, RefWorks:340, RefWorks:339}. Instead, studies have focused 
predominantly on candidates for school board, why they run, and if their races 
are competitive \citep{RefWorks:311, RefWorks:321, RefWorks:322, RefWorks:207}. 
Yet, it might be argued, that such studies are jumping the gun. The conventional 
wisdom about school board elections is that they are minor affairs, settled on 
obscure issues, and conducted among large public apathy. Off-cycle school board 
elections are decidedly low turnout affairs. If boards have competitive elections 
that lead to consequential decisions like superintendent turnover (ala Dissatisfaction 
Theory \citep{RefWorks:297}), what does it matter if those hotly contested electoral 
defeats of incumbents came at the hand of 10\% of the electorate? The overwhelming 
majority has indeed cast a vote of indifference. 

By leaving out the question of turnout, Dissatisfaction Theory has left out an 
important strategic actor in the act of expressing dissatisfaction with the 
direction of the schools -- the voter \citep{RefWorks:243}. The current study 
rectifies this by providing a descriptive analysis of voter turnout over an extended 
number of election cycles at hundreds of school districts within Wisconsin simultaneously. 
Naturally it is not expected to find school board races to be high turnout affairs 
year in and year out. But, if Dissatisfaction Theory is to be saved from criticism 
that it ignores voters in favor of focusing on board-superintendent relationships, 
some cases of heightened turnout must be identified. The current study is fruitful 
ground for examining voter turnout trends in school board elections as well as exploring 
possible opportunities for voters to express dissatisfaction with academic, fiscal, 
or the policy performance of the local school district. 

Several competing conceptions of voter turnout can be explored at the local level 
through the lens of turnout in school board elections. Table \ref{turnoutH} lists 
the evidence in school board turnout that would be consistent with the dominant 
theories of voter turnout. 


\begin{table}[h!]
\begin{center}
\begin{small}
\caption{Theories of Turnout and Expected Evidence}
\label{turnoutH}
\begin{tabular}{| p{.15\textwidth} | p{.2\textwidth} | p{.6\textwidth} |}
\hline
\textbf{Theory}  & \textbf{Citation} & \textbf{Expected Finding} \\ \hline
Habit-forming voters  & \citet{Plutzer2002} & High autocorrelation in turnout\\ \hline
Game-theoretic voters & \citet{Mueller3, aldrich1993} & Smaller electorate has higher turnout \\ \hline
Social group led voting &\citet{Fowler2005} & Union membership as share of electorate drives up turnout \\ \hline
Pivotal-voter & \citet{Downs57} & Small electorate and close election drives turnout up \\ \hline
Social group led voting &\citet{Fowler2005} & Polarized communities have higher turnout \\ \hline
\end{tabular}
\end{small}
\end{center}
\end{table}

\subsection{Exogenous State Level Policy Shocks and Voter Decisions}

There is one specific theory about voter turnout that this data can test is the 
linkage between voter turnout in school board elections and state level policy 
shocks. While there has recently been more attention on the linkage between 
national political trends and state level election returns (CITATION NEEDED), these 
studies have not been extended to the local level. However, the local level provides 
an interesting level of analysis because local elections are much more diverse than 
their state level counterparts in terms of the timing, partisanship of office, and 
structure of the elections. Where it is difficult to disentangle relationships between 
national legislature and state legislature composition from national trends in 
partisanship and strategic redistricting (CITATION NEEDED), the diversity of local 
elections allows more opportunities for exogenous variation to explore the causal 
transmission of voter preferences at the state level being expressed in local 
elections. 

Wisconsin from 2010-2012 provides an ideal test case for the transmission of 
state level political preferences to local non-partisan office. As discussed in 
Chapter \ref{chap:theoryreview} and illustrated in Figure \ref{fig:wiscotimeline} 
the election of Scott Walker to Governor and the subsequent passage of Act 10, the 
2011-13 biennial budget, and the mass protests in Madison, WI, state level politics 
in Wisconsin became suddenly extremely polarized.\footnote{Perhaps the most salient 
issue is the requirement under the law that any wage increase greater than a cost 
of living adjustment pegged to CPI be approved by a district wide referenda - 
effectively placing a firm ceiling on any salary negotiations2011 Wisconsin Act 10: 
https://docs.legis.wisconsin.gov/2011/related/acts/10} In addition to the sudden 
polarization, much of the rhetoric, protest, and contentiousness surrounding 
the Walker administration was focused squarely on education politics and labor 
relations - issues that school boards have great local influence over. In fact, 
Act 10 represented a sweeping expansion of the power of school boards away from 
local teachers' unions. Thus, the composition of the school board mattered more 
than ever after Act 10. Importantly, the passage of Act 10 was completely unforeseen 
by political observers and the public and represent an exogenous shock to the 
state political system. \footnote{This last point is somewhat disputed by some 
political observers and Scott Walker himself, but the consensus has been that 
this was not the case according to the \href{http://www.politifact.com/wisconsin/statements/2011/feb/22/scott-walker/wisconsin-gov-scott-walker-says-he-campaigned-his-/}{Milwaukee Journal Sentinel}.}

This paper leverages the unexpected passage of Act 10 as providing a natural 
experiment of the effect of state level policy shocks on local level electoral 
participation. In fact, in this case, the policy shock and surrounding political 
polarization was directly related to the powers of school boards, a locally 
elected non-partisan office. Even more interestingly, school board elections are 
held off-cycle in Wisconsin in the spring on the $2^{nd}$ Tuesday each April. This 
combination of non-partisan off-cycle elections and a wide and deeply salient policy 
shock provide an ideal test of whether or not state level policy actions can spur 
voters in local elections. Specifically, following \citet{Matsusaka}, the state 
level policy shock should have increased dramatically the information available to 
voters about school boards, school board members, and the impact of Act 10 on their 
community. Thus, turnout should have increased above historical levels. 
If there is no impact on voter turnout in the wake of such a broad expansion of 
school board power over an issue as polarizing as collective bargaining rights 
for local employees, we would expect smaller policy shocks to have no effect either. 
This would represent a rejection of the applicability of information theory to 
explain turnout in low-salience non-partisan elections. This would also provide 
evidence in support of dissatisfaction theory - with heightened stakes for board 
policy preferences.\footnote{One concern might be that the polarization was so 
extreme and the state so saturated with negative campaigning of all forms that 
voter turnout was suppressed even in the case of spring school board elections 
\citep{NegativeCamp1999}.}

\begin{table}[h!]
\begin{center}
\begin{small}
\caption{Theories of Turnout and Expected Evidence: Causal}
\label{turnoutH2}
\begin{tabular}{| p{.15\textwidth} | p{.2\textwidth} | p{.6\textwidth} |}
\hline
\textbf{Theory}  & \textbf{Citation} & \textbf{Expected Finding} \\ \hline
Information theory  & \citet{Matsusaka} & Large turnout increases after passage of Act 10\\ \hline
Dissatisfaction theory & \citet{RefWorks:297,RefWorks:271} & Large turnout increase after Act 10 \\ \hline
\end{tabular}
\end{small}
\end{center}
\end{table}

This follows along the work of \citet{Lassen2005} who leveraged an exogenous shock 
of an information campaign in some electoral districts and found that 
increasing voters' information also produced a sizeable effect on the propensity 
to vote. 

%% Consider partisan preferences revealed by recall petition -- anecdotal evidence 
%% that this influenced voters in school board elections
%% Can a subsample of school districts be chosen to look up recall petition activity 
%% Among incumbents and candidates?


\section{Data}
\label{sec:turnData}

%% Standard data writeup here

As we saw in Chapter \ref{chap:wisco}, the spring elections in Wisconsin are decidedly low turnout affairs. Estimating school board election turnout is decidedly difficult because it requires 
both a reliable assessment of the votes cast in the election and a reliable 
estimate of the voting age population in the school district. Due to the fact that 
school districts have diverse organizational structure above, and no requirements 
on election reporting to a central authority, it is often difficult to ascertain 
the number of votes cast by voters within the school district. To determine the voting 
age population (VAP) of the school district in Wisconsin, I employ the finest resolution 
estimate of voting age population available -- VAP estimates by minor civil division (MCD) 
conducted by the Wisconsin Department of Administration described in Chapter \ref{chap:wisco} Section \ref{subsec:turnout}.\footnote{Available online: \url{http://gab.wi.gov/elections-voting/results}. Statistics for Wisconsin Minor Civil Divisions are maintained by the Department of Administration Demographic Services Center: \url{http://www.doa.state.wi.us/section\_detail.asp?linkcatid=11\&linkid=64\&locid=9}}

With the method described in that section, I am able to estimate the voter turnout for 
each school district in Wisconsin on school board elections as well as on presidential, 
gubernatorial, and non-partisan spring elections. 





\begin{knitrout}
\definecolor{shadecolor}{rgb}{0.969, 0.969, 0.969}\color{fgcolor}\begin{figure}[]

\includegraphics[width=\maxwidth]{../../phd_figs/chap4-turnoutComparison} \caption[Comparison of Fall and Spring Turnout in Wisconsin Elections]{Comparison of Fall and Spring Turnout in Wisconsin Elections\label{fig:turnoutComparison}}
\end{figure}


\end{knitrout}


\begin{knitrout}
\definecolor{shadecolor}{rgb}{0.969, 0.969, 0.969}\color{fgcolor}\begin{figure}[]

\includegraphics[width=\maxwidth]{../../phd_figs/chap4-partyDivision} \caption[Comparison of School Board Turnout to Electorate Polarization]{Comparison of School Board Turnout to Electorate Polarization\label{fig:partyDivision}}
\end{figure}


\end{knitrout}

\begin{knitrout}
\definecolor{shadecolor}{rgb}{0.969, 0.969, 0.969}\color{fgcolor}\begin{figure}[]

\includegraphics[width=\maxwidth]{../../phd_figs/chap4-turnoutCompetitive} \caption[Turnout in School Board Elections based on Contestation]{Turnout in School Board Elections based on Contestation\label{fig:turnoutCompetitive}}
\end{figure}


\end{knitrout}

\begin{knitrout}
\definecolor{shadecolor}{rgb}{0.969, 0.969, 0.969}\color{fgcolor}\begin{figure}[]

\includegraphics[width=\maxwidth]{../../phd_figs/chap4-turnoutSize} \caption[Turnout in School Board Elections based on Electorate Size]{Turnout in School Board Elections based on Electorate Size\label{fig:turnoutSize}}
\end{figure}


\end{knitrout}


\begin{knitrout}
\definecolor{shadecolor}{rgb}{0.969, 0.969, 0.969}\color{fgcolor}\begin{figure}[]

\includegraphics[width=\maxwidth]{../../phd_figs/chap4-turnoutAutocorrelation} \caption[Turnout in School Board Elections based on Prior Turnout]{Turnout in School Board Elections based on Prior Turnout\label{fig:turnoutAutocorrelation}}
\end{figure}


\end{knitrout}


% Turnout and group membership

\begin{knitrout}
\definecolor{shadecolor}{rgb}{0.969, 0.969, 0.969}\color{fgcolor}\begin{figure}[]

\includegraphics[width=\maxwidth]{../../phd_figs/chap4-unionInf} \caption[Turnout in School Board Elections based on School Staff Share]{Turnout in School Board Elections based on School Staff Share\label{fig:unionInf}}
\end{figure}


\end{knitrout}

% Competitiveness

\begin{knitrout}
\definecolor{shadecolor}{rgb}{0.969, 0.969, 0.969}\color{fgcolor}\begin{figure}[]

\includegraphics[width=\maxwidth]{../../phd_figs/chap4-competition} \caption[Competitiveness and Turnout]{Competitiveness and Turnout\label{fig:competition}}
\end{figure}


\end{knitrout}

I begin with the explanations in Table \ref{turnoutH} and explore them using the Wisconsin 
data. I start with the models based on electorate size and explore if at the 
extremely small electorate sizes common in school board elections there is evidence 
that voters are more likely to turnout than in larger school board races. Figure 
\ref{fig:turnoutSize} shows that there is a consistent negative relationship 
between electorate size and school board election turnout across the years of data 
available in Wisconsin. As voting age population increases, Wisconsin school board 
elections have lower turnout with a correlation of 
-0.134. 
This conforms with the predictions of game-theoretic approaches which state that 
in a small electorate a voter can expect their vote to matter more in the outcome and 
thus increase the potential benefits associated with voting
\citep{Mueller3, aldrich1993, Downs57}. However, this does not represent a full test 
of this theory, as voters are also voting up ticket in statewide non-partisan races 
such as Supreme Court and State Superintendent, and school board turnout may be 
driven by these races. 

Having competitive races is critical in driving voter turnout, as voters have 
very little reason to vote if they are not presented with a meaningful choice \citep{Downs57}. 
Figure \ref{turnoutCompetitive} shows voter turnout across the competitiveness typologies 
described in Chapter \ref{chap:wisco} with a clear pattern of higher turnout in races 
featuring an incumbent and a serious challenger. Open seats also appear to feature higher 
turnout than uncontested races. The turnout in uncontested races in the bottom panel 
is very likely being driven by the presence of other contested races on the ticket. 
Using a statistical model we will be able to explore this in greater detail. 


Figure \ref{fig:turnoutAutocorrelation} explores how habit forming turnout may be 
in spring elections in school districts \citep{Plutzer2002}. Using a two-year lag 
of turnout, Figure \ref{fig:turnoutAutocorrelation} shows a consistently high 
positive correlation between lagged and current turnout, averaging 
0.349.
\footnote{For a single year lag, the corresponding correlation is 
0.311 and the 
year to year pattern is largely similar.} A robust linear model is used to show 
the smoothed relationship between lagged and current turnout. Of note is the pattern 
shift in both 2005 and 2011 where the correlation is more pronounced and turnout 
appears to be inflated. Figure \ref{fig:turnoutComparison} looks at this further 
yet by showing the relationship between the average of the turnout in the prior two 
fall elections for Governor and President compared to the current school board turnout. 
These correlate positively as well at a level of 
0.188. 

With a reduced electorate, school board elections are an interesting case to study 
the role substantial interest groups can play in an election. \citet{Fowler2005} 
provides the premise that cohesive social groups can increase turnout by activating 
their members to vote. \citet{RefWorks:207} suggests that unionized teachers are 
just one such group and provides evidence of school board elections being influenced 
substantially by such groups. To evaluate this, Figure \ref{fig:unionInf} shows 
the turnout in the school board election compared to the percentage of the twice lagged 
school board electorate that licensed teachers would make up if they all voted. 
This comparison illustrates how large of a voting block a cohesive teacher union 
might estimate they possess based on the number of votes cast in previous elections. 
The relationship shown in Figure \ref{fig:unionInf} is not what theory would expect. 
Here a higher teacher share of the prior school board turnout is associated with depressed turnout rates, 
instead of elevated turnout. The average correlation across years is 
-0.124. 
I will pursue this finding further with statistical models, 
but it is worth noting that there may be limits on the impact social groups can have 
on driving up turnout before they capture the electorate and thus reduce competitiveness 
and turnout. In fact, for \citet{RefWorks:207} this would represent the ideal condition 
for a union, freeing their members from the need to be politically active and influencing 
the school board through the nomination and campaign process instead. 

Still, when school board races are competitive, do they exhibit much higher turnout? Figure 
\ref{fig:competition} shows that races where the winner and loser are within 15\% of 
one another in votes do exhibit higher turnout.\footnote{For this graph the definition 
of win-margin described in Chapter \ref{chap:wisco} was used.} About 1/3 races is 
competitive in this way. This is an imperfect measure however, because voters do 
not know before they go to the ballot how competitive a race will be. In fact, in 
school board races, it is very unlikely that voters have much information at all about 
how close the race will be due to a lack of polling information and very little 
campaigning common in school board races \citep{RefWorks:321, RefWorks:217, RefWorks:322}. 

Figure \ref{fig:partyDivision} takes a different approach to this problem by looking 
at the information voters may have about their community -- the level of partisan 
polarization. Using the average Democratic share of the two-party vote from the prior 
Gubernatorial and Presidential election, I explore the variation in school board election 
turnout. Theory would suggest that the more divided a community is on partisan issues, the 
greater the impetus to vote \citep{Fowler2005}. Thus in Figure \ref{fig:partyDivision} we 
should expect to see an inverted U shape representing higher turnout around 50\% share of 
the two party vote. In fact, this pattern does not occur in any year, and in some years 
there is evidence of just the opposite relationship -- a correlation between higher 
partisan unity and turnout. Figure \ref{fig:partisan2} shows the linear relationship 
between higher partisan unity and turnout. In most years there is almost no relationship, 
but in 2011 and 2012 there appears to be an uptick in school board turnout associated 
with more strongly partisan school districts. The correlation is 
0.134. 

\begin{knitrout}
\definecolor{shadecolor}{rgb}{0.969, 0.969, 0.969}\color{fgcolor}\begin{figure}[]

\includegraphics[width=\maxwidth]{../../phd_figs/chap4-partisan2} \caption[Another Look at Relationship Between Turnout and Partisanship]{Another Look at Relationship Between Turnout and Partisanship\label{fig:partisan2}}
\end{figure}


\end{knitrout}

Taken alone, these figures do not suggest that any of the theories in Table \ref{turnoutH} 
are confirmed or denied, but it does help illustrate the diversity of school board 
elections in Wisconsin and the need for statistical models to help separate out 
competing explanations that may themselves by highly correlated or confound 
one another. 

\clearpage

\subsection{Control Variables}

In addition to the variables of interest shown above, it is necessary to include 
some common demographic control variables. Political scientists have long established 
the strong relationship between certain demographic characteristics and voter 
turnout in partisan elections. % CITENEEDED
In order to understand the effect of other variables on turnout, it is desirable to 
remove as many other explanations as possible given the availability of data and 
the ability to generate valid estimates. Fortunately, using US census data, an 
abundance of data is available on the population of Wisconsin school districts. 

% Describe per_white_all, per65o, PerBachelorOrAbove, OOH_share, PCI
% also possibly millrate?



\clearpage
\section{Methods}

This paper has two main methodological approaches to evaluating theories of voter 
turnout in school board elections. The first follows on several empirical studies 
of theories of voter turnout to compare the power of various theories of voter 
turnout to one another \citep{Larcinese, Lapp99, KanYang2001,Coate2008582}. The 
second approach leverages the exogenous shock of the Walker administration in 
Wisconsin and the passage of Act 10 to evaluate the causal impact of changes in 
voter information on both the level of turnout in school board elections, and 
the prompting of voters to exercise their dissatisfaction through the ballot. 

A key factor to control for in the case of school boards is the fiscal issue dimension. 
Researchers have studied the impact of economic conditions on voting before, and 
have found some evidence that economic conditions influence voting patterns 
\citep{EconandOutcomes2000}. 
If we reduce school boards to a two-dimensional game with incumbents and non-incumbents 
along one dimension and the relative fiscal performance along the other, ceteris 
paribus we would expect that incumbents are at an advantage (and turnout is lower) 
when fiscal performance is strong. To measure fiscal performance we use the share 
of the revenue limit a district is spending.

\subsection{Evaluating Theories in Non-Partisan Local Elections}

Now we turn to operationalizing Table \ref{turnoutH} into estimable equations to 
look for evidence of voter turnout theories in school board elections. 

\[(1) \hspace{0.4 in} VP_{t} = \alpha + \beta X_{lag} + \psi Ce_{t}+\epsilon \] \\

Here $VP_{t}$ is voter turnout in time $t$. $\alpha$ is the intercept and $x_{lag}$ 
represents a vector of lagged covariates such as turnout in school district $i$. 
Additionally, we include a term $Ce_{t}$ for the level of candidate participation 
in the election as a control. As per \citet{habit02}, a high predictive value 
of the lagged term, especially in smaller districts, would suggest the electorate 
is largely consistent of the same individuals year to year, suggesting that learning 
is occurring. To test this more rigorously, we could specifically look at years 
where turnout increases are particularly sharp, and see if these correspond to 
heightened turnout in particular years through the inclusion of a $\delta turnout$ 
term. 

Next, we turn to testing the pivotal-voter model using a similar method as \citep{Coate2008582}. 
The pivotal-voter model predicts that turnout should be higher in elections that 
are closer. % OR IS IT THE OPPOSITE? 
Thus we include the $M_{t}$ for the margin of victory for candidates in the school 
board election in the equation below. 

%% do we consider controlling for the effect of candidate participation?
%% Consider including candidate participation in table 4.1?

\[(2) \hspace{0.4 in}  VP_{t} = \alpha + \beta X_{lag} + \theta M_{t} + \psi Ce_{t} + \epsilon \] \\

The size of the electorate is important in many theories of voter turnout. \citet{dahl1973size} 
first hypothesized that the motivation to participate would be higher in small 
electorates, and there has been some evidence to support this \citep{Frandsen2002}. 
Thus, in equation three we include a $\theta S_{t}$ term that controls for the 
size of the electorate. In a game-theoretic perspective, the smaller electorate 
would drive turnout because the costs of non-voting can be more easily estimated 
by voters in small electorates, and the expected benefit of such non-voting is 
larger. 

\[(3) \hspace{0.4 in}  VP_{t} = \alpha + \beta X_{lag} + \theta S_{t} + \psi Ce_{t} + \epsilon \] \\

A pure \citet{Downs57} model would suggest that as the value of PB increases, then 
so should turnout. Thus, equation 4 contains terms for both the size of the electorate 
and the closeness of victory $ \theta S_{t} *\theta M_{t}$. Unfortunately, this 
equation is not able to estimate the value of B, or the expected benefit of 
the candidate in question to a voter. This will be addressed in the next section. 

\[(4) \hspace{0.4 in}  VP_{t} = \alpha + \beta X_{lag} + \theta S_{t} *\theta M_{t}  + \epsilon \] \\

Finally, following \citet{Fowler2005}, we investigate the impact that social group 
membership may have on voting. In equation 5, $P_{t}$ represents a measure of the
polarization in the community. There are two options for measuring $P_{t}$ - either 
by looking at the partisan makeup of the electorate using the partisanship of the 
school district in partisan statewide elections, or looking at teacher membership 
as a share of the total electorate. There has been little work explicitly testing the 
former, but some evidence that the share of teachers in the electorate is an 
important driver of turnout and policy change in school board elections \citep{RefWorks:207, 
RefWorks:218,RefWorks:320}. 

\[(5) \hspace{0.4 in}  VP_{t} = \alpha + \beta X_{lag} + \theta S_{t} 
+ \gamma P_{t} + \psi Ce_{t} + \epsilon \] \\

\subsection{The Impact of Salience and Information on Voter Turnout}

While the equations above can prove useful for providing evidence of various theories 
of voter turnout in non-partisan local elections, they cannot test some theories and 
they do not shed any light on causes of voter turnout. Table \ref{turnoutH2} demonstrates 
how two such models can be explicitly tested by leveraging the exogenous shock of 
the passage of Act 10. To estimate these effects, we estimate an equation like below: \\

\[(6) \hspace{0.4 in} VP_{t} = \alpha + \beta X_{lag} + \gamma W_{t} + \lambda P_{lag} +
      \tau 1- |W_{t} * P_{lag}| + \theta I_{lag} + \psi Ce_{t}+\epsilon \] \\

Here $VP_{t}$ is voter turnout in either the 2011 or 2012 Spring Elections by school 
district in Wisconsin. $\alpha$, $X_{lag}$ remain the same as above though now 
including some measure of prior turnout in the district such as turnout in the 
Gubernatorial election in 2010.  $Z_{t}$ is the absolute value of the deviation 
from 50\% of the vote share for the Governor. $I_{lag}$ representing an indicator 
of the strength of teachers' unions in the district--likely the most salient and 
relevant interest group in the spring elections \citep{RefWorks:320, RefWorks:207, RefWorks:218}.
Voters are unlikely to participate if presented with no choice, so $Ce_{t}$ represents 
a measure of the number of candidates per seat. $W_{t}$ represents the vote share 
for Governor Walker in the recall election--a measure of community support or opposition 
to the Governor's education policy. Thus it would be expected that $\theta$ should 
be positive--strong interest groups should drive political activity. The inclusion 
of partisanship predictors of voter turnout could also be explored here. The 
effects of $\gamma$  and $\lambda$ alone are uninteresting, but the interaction 
of $W_{t} * P_{lag}$ is of great interest. This essentially represents a measure 
of board policy alignment with voter preferences--taking the absolute value simplifies 
interpretation. As such, it should be expected that for large values of $W_{t} * P_{lag}$, 
voter participation would be small as the board policy and the community 
preferences would be aligned. As $W_{t} * P_{lag}$ approaches zero, voter 
participation will increase as predicted by \citep{RefWorks:243}--voters seeing 
dissonance between their preferences and board policy are more likely to vote. 
Thus $\tau$ should be negative.  %Add detail here

\subsection{Model Specifications}

Voter turnout is driven by many forces. Aside from the exogenous shock provided 
by Scott Walker's election and the passage of Act 10, the ability to test the 
causal impact of any factor on school board elections is limited. In my first 
set of models I look for evidence that the explanations described in Table \ref{turnoutH} 
are supported in the data available on Wisconsin school boards. In each case, 
I start with a specification that pools observations across years and districts and 
estimates the impact of the key independent variable and a few key demographic 
control variables described in Section \ref{sec:turnData}. 

Due to the presence of outliers exhibited in the figures in \ref{sec:turnData} my 
initial models will represent penalized regressions estimated using the \texttt{ols} 
function in the \texttt{rms} package for R \citep{cran, rms}. Instead of estimating 
$\beta$ as $(X'X)^{-1} X'Y$, a user specified penalty factor P is included 
to adjust $\beta$ and estimate it as $(X'X + P)^{-1} X'Y$. 

%Demonstrate outliers
% Talk about penalized maximum likelhood
% Talk about the problem of overfit with fixed effects in large amount
% Talk about mixed effect regression 
% Bayesian and tobit regression options
% Triangulating 


\section{Results}

\subsection{Autocorrelated Turnout}

The initial pooled linear regression is shown in Table \ref{tab:lagmods} as the 
first column. This model includes no adjustments for clustering to the standard 
errors. It is included here to demonstrate the need for specification testing 
in the Wisconsin dataset due to the ragged panel structure of the data. 

\begin{knitrout}
\definecolor{shadecolor}{rgb}{0.969, 0.969, 0.969}\color{fgcolor}\begin{figure}[]

\includegraphics[width=\maxwidth]{../../phd_figs/chap4-simpleModel} \caption[Regression Diagnostics for A Simple Multiple Regression]{Regression Diagnostics for A Simple Multiple Regression\label{fig:simpleModel}}
\end{figure}


\end{knitrout}

Figure \ref{fig:simpleModel} shows regression diagnostics for this simple multiple 
regression. The key result here is that there are a number of high leverage 
observations that appear to be extreme outliers. Additionally, the upper left 
quadrant of Figure \ref{fig:simpleModel} suggests the presence of some level of 
heteroskedacticity in the results, a likely result of the clustering and repeated 
observation features of the dataset. 

\begin{table}[!ht]
\caption{Models of Lagged Turnout in School Board Elections}
\label{tab:lagmods} 
\begin{tabular}{ l D{.}{.}{2}D{.}{.}{2}D{.}{.}{2}D{.}{.}{2} } 
\hline 
  & \multicolumn{ 1 }{ c }{ Naive OLS } & \multicolumn{ 1 }{ c }{ Simple RLM } & \multicolumn{ 1 }{ c }{ RLM + Year FE } & \multicolumn{ 1 }{ c }{ RLM + District FE } \\ \hline
 %                  & Naive OLS         & Simple RLM        & RLM + Year FE     & RLM + District FE\\ 
(Intercept)        & -72.08 ^*         &                   &                   &                  \\ 
                   & (14.11)           &                   &                   &                  \\ 
turnoutLag2        & 26.30 ^*          & 26.30 ^*          & 27.06 ^*          & -1.32            \\ 
                   & (2.13)            & (4.67)            & (4.59)            & (4.46)           \\ 
fallTurnout        & 12.00 ^*          & 12.00 ^*          & 15.83 ^*          & 10.95            \\ 
                   & (3.01)            & (4.82)            & (4.93)            & (6.72)           \\ 
VAP\_adj          & -2.04 ^*          & -2.04 ^*          & -2.02 ^*          & -23.92 ^*        \\ 
                   & (0.21)            & (0.29)            & (0.28)            & (6.64)           \\ 
per\_white\_all  & -5.39 ^*          & -5.39             & -3.24             & 28.98            \\ 
                   & (2.64)            & (3.66)            & (3.87)            & (40.77)          \\ 
per65o             & 29.04 ^*          & 29.04 ^*          & 27.25 ^*          & 107.16 ^*        \\ 
                   & (6.14)            & (10.88)           & (10.87)           & (46.88)          \\ 
PerBachelorOrAbove & -8.66             & -8.66             & 0.47              & -17.42           \\ 
                   & (5.64)            & (7.26)            & (7.21)            & (27.48)          \\ 
OOH\_share        & -2.37             & -2.37             & -1.66             & 15.69            \\ 
                   & (2.13)            & (3.18)            & (3.08)            & (19.13)          \\ 
PCI                & 9.79 ^*           & 9.79 ^*           & 3.64 ^*           & 25.13 ^*         \\ 
                   & (1.48)            & (1.81)            & (1.73)            & (3.38)           \\ 
Intercept          &                   & -72.08 ^*         & -17.18            & -77.79           \\ 
                   &                   & (16.95)           & (16.79)           & (82.75)           \\
 $N$                & 2390              & 2390              & 2390              & 2390             \\ 
$R^2$              & 0.22              & 0.22              & 0.36              & 0.43             \\ 
adj. $R^2$         & 0.22              & 0.22              & 0.35              & 0.35             \\ 
Resid. sd          & 8.99              & 8.99              & 8.18              & 8.20              \\ \hline
 \multicolumn{5}{l}{\footnotesize{Standard errors clustered at the school district except in model 1.}}\\
\multicolumn{5}{l}{\footnotesize{$^*$ indicates significance at $p< 0.05 $}}\\
\multicolumn{5}{l}{\footnotesize{Year and district fixed effects omitted.}} 
\end{tabular} 
 \end{table}


Turnout appears to be very highly correlated as shown in Table \ref{tab:lagmods}. 
Model 1 represents a standard pooled regression with clustered standard errors. 
The key independent variable is twice lagged turnout for school board elections, which 
represents our test of habit forming in school board elections. Control variables 
included are some standard demographic predictors of voter turnout including per 
capita income, education level, percent of population over 65, percent of population 
that is white. Additionally, a school district specific control, the share of residents 
in owner-occupied housing is included to reflect the tight coupling between school boards 
and property owning. 

Using twice-lagged school board turnout. Results are not substantively different 
when I estimate using single lagged school board turnout, as shown in Table \ref{tab:lagmods2}. 





We should also fit some better models. Partial pooling models can be fit in R 
using the \texttt{lme4} package \citep{lme4, gelhill}. 

\begin{knitrout}
\definecolor{shadecolor}{rgb}{0.969, 0.969, 0.969}\color{fgcolor}\begin{kframe}
\begin{verbatim}
## Linear mixed model fit by REML ['lmerMod']
## Formula: 
## turnout ~ turnoutLag2 + fallTurnout + log(VAP_adj) + per_white_all +  
##     per65o + PerBachelorOrAbove + OOH_share + log(PCI) + (1 |      distid)
##    Data: moddat
## 
## REML criterion at convergence: 17155
## 
## Scaled residuals: 
##    Min     1Q Median     3Q    Max 
## -4.390 -0.597 -0.114  0.489  7.923 
## 
## Random effects:
##  Groups   Name        Variance Std.Dev.
##  distid   (Intercept) 12.8     3.58    
##  Residual             69.9     8.36    
## Number of obs: 2390, groups:  distid, 309
## 
## Fixed effects:
##                    Estimate Std. Error t value
## (Intercept)         -253.88      43.67   -5.81
## turnoutLag2           11.86       2.20    5.39
## fallTurnout           10.93       3.95    2.76
## log(VAP_adj)         -23.86       2.57   -9.30
## per_white_all         -7.09       3.75   -1.89
## per65o                39.00       8.73    4.47
## PerBachelorOrAbove   -10.98       7.44   -1.48
## OOH_share             -1.40       3.03   -0.46
## log(PCI)             137.11      19.26    7.12
## 
## Correlation of Fixed Effects:
##             (Intr) trntL2 fllTrn l(VAP_ pr_wh_ per65o PrBcOA OOH_sh
## turnoutLag2 -0.079                                                 
## fallTurnout  0.153 -0.058                                          
## log(VAP_dj) -0.231  0.231  0.073                                   
## per_whit_ll  0.005  0.001 -0.132  0.057                            
## per65o       0.177 -0.079 -0.048  0.116 -0.226                     
## PrBchlrOrAb  0.629 -0.062 -0.432 -0.274  0.158  0.243              
## OOH_share    0.194  0.019 -0.130  0.117 -0.356  0.628  0.193       
## log(PCI)    -0.986  0.046 -0.187  0.089 -0.065 -0.230 -0.609 -0.238
\end{verbatim}
\end{kframe}
\end{knitrout}

Problems abound. 

\subsection{Electorate Size and Margin of Victory}

\begin{table}[!ht]
\caption{Models of Competitiveness and Turnout in School Board Elections}
\label{tab:closemods} 
\begin{tabular}{ l D{.}{.}{2}D{.}{.}{2}D{.}{.}{2} } 
\hline 
  & \multicolumn{ 1 }{ c }{ Simple RLM } & \multicolumn{ 1 }{ c }{ RLM + Year FE } & \multicolumn{ 1 }{ c }{ RLM + District FE } \\ \hline
 %                     & Simple RLM        & RLM + Year FE     & RLM + District FE\\ 
Intercept             & -0.68 ^*          & -0.13             & -0.58            \\ 
                      & (0.20)            & (0.21)            & (0.80)           \\ 
CLOSE=Not Competitive & -0.05 ^*          & -0.05 ^*          & -0.05 ^*         \\ 
                      & (0.00)            & (0.00)            & (0.00)           \\ 
VAP\_adj             & -0.03 ^*          & -0.03 ^*          & -0.24 ^*         \\ 
                      & (0.00)            & (0.00)            & (0.07)           \\ 
per\_white\_all     & -0.04             & -0.04             & 0.06             \\ 
                      & (0.04)            & (0.05)            & (0.38)           \\ 
per65o                & 0.38 ^*           & 0.35 ^*           & 1.15 ^*          \\ 
                      & (0.15)            & (0.14)            & (0.44)           \\ 
PerBachelorOrAbove    & 0.08              & 0.04              & -0.12            \\ 
                      & (0.08)            & (0.09)            & (0.26)           \\ 
OOH\_share           & -0.02             & -0.03             & 0.19             \\ 
                      & (0.04)            & (0.04)            & (0.18)           \\ 
PCI                   & 0.12 ^*           & 0.05 ^*           & 0.25 ^*          \\ 
                      & (0.02)            & (0.02)            & (0.03)           \\ 
fallTurnout           &                   & 0.18 ^*           & 0.12             \\ 
                      &                   & (0.06)            & (0.07)           \\ 
turnoutLag2           &                   &                   & 0.00             \\ 
                      &                   &                   & (0.04)            \\
 $N$                   & 2390              & 2390              & 2390             \\ 
$R^2$                 & 0.21              & 0.36              & 0.47             \\ 
adj. $R^2$            & 0.21              & 0.36              & 0.39             \\ 
Resid. sd             & 0.09              & 0.08              & 0.08              \\ \hline
 \multicolumn{4}{l}{\footnotesize{Robust standard errors in parentheses}}\\
\multicolumn{4}{l}{\footnotesize{$^*$ indicates significance at $p< 0.05 $}}\\
\multicolumn{4}{l}{\footnotesize{Year and district fixed effects omitted.}} 
\end{tabular} 
 \end{table}


Competitiveness does matter. 

\begin{knitrout}
\definecolor{shadecolor}{rgb}{0.969, 0.969, 0.969}\color{fgcolor}\begin{kframe}
\begin{verbatim}
## Linear mixed model fit by REML ['lmerMod']
## Formula: turnout ~ CLOSE + turnoutLag2 + fallTurnout + log(VAP_adj) +  
##     per_white_all + per65o + PerBachelorOrAbove + OOH_share +  
##     log(PCI) + (1 | distid)
##    Data: moddat
## 
## REML criterion at convergence: -4899
## 
## Scaled residuals: 
##    Min     1Q Median     3Q    Max 
## -4.435 -0.597 -0.112  0.486  8.268 
## 
## Random effects:
##  Groups   Name        Variance Std.Dev.
##  distid   (Intercept) 0.00130  0.0360  
##  Residual             0.00655  0.0809  
## Number of obs: 2390, groups:  distid, 309
## 
## Fixed effects:
##                      Estimate Std. Error t value
## (Intercept)          -1.04153    0.17644   -5.90
## CLOSENot Competitive -0.04740    0.00377  -12.59
## turnoutLag2           0.12333    0.02133    5.78
## fallTurnout           0.10131    0.03887    2.61
## log(VAP_adj)         -0.03126    0.00294  -10.62
## per_white_all        -0.07197    0.03715   -1.94
## per65o                0.37553    0.08677    4.33
## PerBachelorOrAbove   -0.11824    0.07383   -1.60
## OOH_share            -0.01765    0.03014   -0.59
## log(PCI)              0.14550    0.01845    7.89
## 
## Correlation of Fixed Effects:
##             (Intr) CLOSEC trntL2 fllTrn l(VAP_ pr_wh_ per65o PrBcOA OOH_sh
## CLOSENtCmpt  0.007                                                        
## turnoutLag2 -0.087 -0.021                                                 
## fallTurnout  0.112  0.015 -0.057                                          
## log(VAP_dj) -0.274  0.147  0.212  0.081                                   
## per_whit_ll -0.089 -0.004  0.004 -0.131  0.075                            
## per65o       0.124  0.007 -0.075 -0.045  0.133 -0.223                     
## PrBchlrOrAb  0.624 -0.002 -0.061 -0.432 -0.279  0.143  0.241              
## OOH_share    0.156  0.004  0.019 -0.127  0.128 -0.354  0.630  0.192       
## log(PCI)    -0.959 -0.041  0.047 -0.182  0.084 -0.050 -0.237 -0.611 -0.244
\end{verbatim}
\end{kframe}
\end{knitrout}

\subsection{Unionization}

\subsection{Exogenous Shocks}

\section{Conclusion}

An interesting proposition here is considering whether or not voters in low-salience 
and low-turnout school board elections are a subset of the voters in higher salience 
and higher turnout elections \citep{Cho2008}

\section{Appendix}

Sensitivity checks and further table analysis. 

Single lagged turnout:

\begin{table}[!ht]
\caption{Models of Lagged Turnout (1x) in School Board Elections}
\label{tab:lagmods2} 
\begin{tabular}{ l D{.}{.}{2}D{.}{.}{2}D{.}{.}{2}D{.}{.}{2} } 
\hline 
  & \multicolumn{ 1 }{ c }{ Naive OLS } & \multicolumn{ 1 }{ c }{ Simple RLM } & \multicolumn{ 1 }{ c }{ RLM + Year FE } & \multicolumn{ 1 }{ c }{ RLM + District FE } \\ \hline
 %                  & Naive OLS         & Simple RLM        & RLM + Year FE     & RLM + District FE\\ 
(Intercept)        & -58.64 ^*         &                   &                   &                  \\ 
                   & (14.16)           &                   &                   &                  \\ 
turnoutLag1        & 18.04 ^*          & 17.30 ^*          & 26.49 ^*          & -14.71 ^*        \\ 
                   & (1.93)            & (2.75)            & (2.94)            & (3.25)           \\ 
fallTurnout        & 10.95 ^*          & 10.97 ^*          & 13.47 ^*          & 10.77            \\ 
                   & (3.04)            & (5.23)            & (4.91)            & (7.62)           \\ 
VAP\_adj          & -2.27 ^*          & -21.14 ^*         & -18.70 ^*         & -201.63 ^*       \\ 
                   & (0.21)            & (2.85)            & (2.56)            & (65.64)          \\ 
per\_white\_all  & -5.22             & -5.06             & -2.88             & 29.89            \\ 
                   & (2.68)            & (3.83)            & (3.58)            & (42.22)          \\ 
per65o             & 32.23 ^*          & 33.47 ^*          & 28.51 ^*          & 116.59 ^*        \\ 
                   & (6.16)            & (11.77)           & (10.69)           & (50.66)          \\ 
PerBachelorOrAbove & -4.15             & -3.70             & 2.14              & -14.92           \\ 
                   & (5.64)            & (7.12)            & (6.49)            & (30.27)          \\ 
OOH\_share        & -1.97             & -1.73             & -0.99             & -0.51            \\ 
                   & (2.13)            & (3.34)            & (2.94)            & (21.08)          \\ 
PCI                & 8.70 ^*           & 89.94 ^*          & 36.56 ^*          & 258.80 ^*        \\ 
                   & (1.49)            & (17.95)           & (16.87)           & (37.84)          \\ 
Intercept          &                   & -153.44 ^*        & -43.24            & -182.73          \\ 
                   &                   & (40.12)           & (38.49)           & (165.33)          \\
 $N$                & 2483              & 2483              & 2483              & 2483             \\ 
$R^2$              & 0.20              & 0.20              & 0.36              & 0.44             \\ 
adj. $R^2$         & 0.20              & 0.20              & 0.35              & 0.36             \\ 
Resid. sd          & 9.18              & 9.15              & 8.23              & 8.16              \\ \hline
 \multicolumn{5}{l}{\footnotesize{Standard errors clustered at the school district.}}\\
\multicolumn{5}{l}{\footnotesize{$^*$ indicates significance at $p< 0.05 $}}\\
\multicolumn{5}{l}{\footnotesize{Year and district fixed effects omitted.}} 
\end{tabular} 
 \end{table}

